% concatenated from main_elsevier.tex
%% ============================================================
%%  Discrete Mathematics — Elsevier  (elsarticle template)
%%  Submission version: preprint
%% ============================================================
\documentclass[preprint,12pt]{elsarticle}

%% ---------------------------------------------------------------
%% Packages
%% ---------------------------------------------------------------
\usepackage{amsmath,amssymb,amsthm}
\usepackage{mathtools}
\usepackage{microtype}
\usepackage{xcolor}
\usepackage{enumitem}

\usepackage[colorlinks=true,
            linkcolor=blue!60!black,
            citecolor=blue!60!black,
            urlcolor=blue!60!black]{hyperref}

%% ---------------------------------------------------------------
%% Journal declaration
%% ---------------------------------------------------------------
\journal{Discrete Mathematics}

%% ---------------------------------------------------------------
%% Theorem environments
%% ---------------------------------------------------------------
\theoremstyle{plain}
\newtheorem{theorem}{Theorem}[section]
\newtheorem{lemma}[theorem]{Lemma}
\newtheorem{proposition}[theorem]{Proposition}
\newtheorem{corollary}[theorem]{Corollary}
\newtheorem{conjecture}[theorem]{Conjecture}

\theoremstyle{definition}
\newtheorem{definition}[theorem]{Definition}
\newtheorem{example}[theorem]{Example}
\newtheorem{remark}[theorem]{Remark}

%% ---------------------------------------------------------------
%% Notation macros
%% ---------------------------------------------------------------
\newcommand{\R}{\mathbb{R}}
\newcommand{\N}{\mathbb{N}}
\newcommand{\lam}{\lambda}
\newcommand{\Lam}{\Lambda}
\newcommand{\om}{\omega}
\newcommand{\eps}{\varepsilon}
\newcommand{\splus}{s^{+}}
\newcommand{\sminus}{s^{-}}
\newcommand{\sprad}{\rho}
\newcommand{\tr}{\operatorname{tr}}
\newcommand{\diag}{\operatorname{diag}}
\newcommand{\floor}[1]{\lfloor #1 \rfloor}
\newcommand{\ceil}[1]{\lceil #1 \rceil}
\DeclareMathOperator{\spec}{spec}
\newcommand{\norm}[1]{\|#1\|}
\newcommand{\dedit}{d_{\mathrm{edit}}}
\newcommand{\Kite}{K_{a,b,c}}
\newcommand{\Tur}{T(n,r)}

%% ============================================================
\begin{document}

%% ---------------------------------------------------------------
%%  Front matter
%% ---------------------------------------------------------------
\begin{frontmatter}

\title{The Bollob\'{a}s--Nikiforov Conjecture for Complete Multipartite
       Graphs and Dense $K_4$-Free Graphs}

\author[inst1]{[Piero Giacomelli]\corref{cor1}}
\ead{pgiacome@gmail.com}
\cortext[cor1]{Corresponding author.}
\address[inst1]{IT Department, TENAX GROUP SPA, Verona, Italy}

\begin{abstract}
The Bollob\'as--Nikiforov conjecture asserts that for any graph
$G \neq K_n$ with $m$ edges and clique number $\omega(G)$,
\[
  \lambda_1^2(G) + \lambda_2^2(G)
  \;\leq\;
  2\!\left(1 - \frac{1}{\omega(G)}\right)m,
\]
where $\lambda_1(G) \geq \lambda_2(G) \geq \cdots \geq \lambda_n(G)$
are the adjacency eigenvalues of $G$.
We prove the conjecture for all complete multipartite graphs
$K_{n_1,\ldots,n_r}$ with $n_1 + \cdots + n_r > r$.
The proof computes the full spectrum via a secular equation, establishes
that $\lambda_2 = 0$ whenever the graph has more vertices than parts,
and then applies Nikiforov's spectral Tur\'an theorem.
For $r = 2$ equality holds for every complete bipartite graph $K_{a,b}$
with $a + b \geq 3$; for $r \geq 3$ equality holds if and only if all
parts have equal size.
We also collect and unify the known partial results for $K_4$-free graphs:
the conjecture holds whenever $\chi(G) \leq 3$ (via Lin--Ning--Wu and
Bollob\'as--Nikiforov for weakly perfect graphs), and for all regular
$K_4$-free graphs (via Zhang's theorem).
A stability result for dense $K_4$-free graphs shows that if the spectral
radius is near the Tur\'an maximum then the graph is structurally close
to the balanced complete tripartite graph, and as a consequence the
conjecture holds for all $K_4$-free graphs with $m = \Omega(n^2)$
when $n$ is sufficiently large.
Finally, we identify the precise obstruction preventing a Hoffman-bound
approach from settling the conjecture for $K_4$-free graphs with
independence number $\alpha(G) \geq n/3$.
\end{abstract}

\begin{keyword}
Spectral graph theory \sep Bollob\'as--Nikiforov conjecture \sep
$K_4$-free graphs \sep adjacency eigenvalues \sep
complete multipartite graphs \sep Tur\'an-type problems
\MSC[2020] 05C50 \sep 05C35 \sep 05C15
\end{keyword}

\end{frontmatter}

%% ============================================================
\section{Introduction}
\label{sec:intro}
%% ============================================================

A central theme in spectral graph theory is to bound combinations of
eigenvalues in terms of classical combinatorial parameters.
Nosal~\cite{Nosal1970} proved that the spectral radius satisfies
$\lambda_1(G) \leq \sqrt{m}$ for triangle-free graphs, with equality
if and only if $G = K_{n/2,n/2}$.
Nikiforov~\cite{Nikiforov2002} extended this to all graphs: for any
graph $G$ with $m$ edges and clique number $\omega(G) \geq 2$,
\begin{equation}
  \label{eq:nikiforov}
  \lambda_1(G)
  \;\leq\;
  \sqrt{2\!\left(1 - \frac{1}{\omega(G)}\right)m},
\end{equation}
with equality if and only if $G$ is a balanced complete
$\omega(G)$-partite graph on $n$ vertices with $\omega(G) \mid n$.
Inequality~\eqref{eq:nikiforov} is the \emph{spectral Tur\'an theorem},
since the extremal graph is the Tur\'an graph $T(n,\omega(G))$.

\paragraph{The Bollob\'as--Nikiforov conjecture.}
Bollob\'as and Nikiforov~\cite{BollobasNikiforov2007} proposed
strengthening \eqref{eq:nikiforov} by simultaneously bounding
$\lambda_1^2 + \lambda_2^2$.

\begin{conjecture}[Bollob\'{a}s--Nikiforov~\cite{BollobasNikiforov2007}]
  \label{conj:BN}
  For any graph $G \neq K_n$ with $m$ edges and clique number
  $\omega(G) \geq 2$,
  \[
    \lambda_1^2(G) + \lambda_2^2(G)
    \;\leq\;
    2\!\left(1 - \frac{1}{\omega(G)}\right)m.
  \]
  The equality cases are:
  \begin{enumerate}[label=(\alph*)]
    \item $\omega(G) = 2$: equality holds for every complete bipartite
          graph $K_{a,b}$ with $a + b \geq 3$.
    \item $\omega(G) \geq 3$: equality holds if and only if $G = T(n,\omega(G))$
          is the balanced complete $\omega(G)$-partite graph with all parts of
          equal size.
  \end{enumerate}
\end{conjecture}

The exclusion of $G = K_n$ is necessary.  For the complete graph $K_n$,
one has $\lambda_1 = n-1$ and $\lambda_2 = -1$, giving
$\lambda_1^2 + \lambda_2^2 = (n-1)^2 + 1 > (n-1)^2 = 2(1-1/n)\binom{n}{2}\cdot 2/n$;
a direct computation confirms that $K_n$ violates the bound.
The conjecture is sharp: the balanced complete $r$-partite graph
$T(n,r)$ with $r \mid n$ satisfies $\lambda_2 = 0$
(proved in Theorem~\ref{thm:multipartite} below) and
$\lambda_1^2 = 2(1-1/r)m$, so equality holds throughout.

\paragraph{Known cases.}
Several special cases of Conjecture~\ref{conj:BN} have been established.
Lin, Ning, and Wu~\cite{LinNingWu2021} confirmed the conjecture for all
triangle-free graphs ($\omega = 2$), with equality exactly at $K_{n/2,n/2}$.
Bollob\'as and Nikiforov~\cite{BollobasNikiforov2007} proved it for all
weakly perfect graphs (graphs satisfying $\chi(G) = \omega(G)$), which
in particular settles the $K_4$-free case whenever $\chi(G) \leq 3$.
Zhang (see~\cite{ElphickWocjan2014}) established the conjecture for
all regular graphs.
Kumar and Pragada~\cite{KumarPragada2024} recently proved it for graphs
containing at most $O(m^{3/2-\varepsilon})$ triangles for any fixed
$\varepsilon > 0$.
Liu and Bu~\cite{LiuBu2025} showed the conjecture holds asymptotically
almost surely for Erd\H{o}s--R\'enyi random graphs.

\paragraph{The open case: dense $K_4$-free graphs.}
A $K_4$-free graph can contain as many as $\Theta(m^{3/2})$ triangles:
the balanced tripartite graph $K_{n/3,n/3,n/3}$ has $\Theta(n^3)$
triangles and $m = \Theta(n^2)$ edges.
Consequently the result of Kumar and Pragada does not apply to $K_4$-free
graphs with $\chi(G) \geq 4$, and this is the principal remaining open case.

\paragraph{Our results.}
We establish the following new results.

\begin{theorem}[Complete multipartite graphs]
  \label{thm:multipartite}
  Let $G = K_{n_1,\ldots,n_r}$ with $r \geq 2$,
  $n = n_1 + \cdots + n_r \geq r+1$ (so $G \neq K_r$), and $m$ edges.
  Then
  \[
    \lambda_1^2(G) + \lambda_2^2(G)
    \;\leq\;
    2\!\left(1 - \frac{1}{r}\right)m,
  \]
  with the following equality characterization:
  \begin{enumerate}[label=(\alph*)]
    \item If $r = 2$: equality holds for \emph{every} complete bipartite graph
          $K_{n_1,n_2}$ with $n_1 + n_2 \geq 3$.
    \item If $r \geq 3$: equality holds if and only if $n_1 = \cdots = n_r$
          (i.e., $G = T(n,r)$ is the balanced complete $r$-partite graph).
  \end{enumerate}
\end{theorem}

\begin{remark}
  Theorem~\ref{thm:multipartite} also follows from the result of
  Bollob\'as and Nikiforov~\cite{BollobasNikiforov2007} for weakly perfect
  graphs (graphs satisfying $\chi(G) = \omega(G)$), since every complete
  $r$-partite graph satisfies $\chi = \omega = r$.
  Our contribution is a self-contained proof via the secular equation
  that yields the exact eigenvalue structure of $K_{n_1,\ldots,n_r}$
  and gives the precise equality characterization above.
  In particular, the equality statement for $r = 2$ (every complete
  bipartite graph achieves equality) does not appear to be recorded
  explicitly in the literature.
\end{remark}

\begin{theorem}[$K_4$-free graphs with $\chi \leq 3$]
  \label{thm:3colorable}
  Let $G$ be a $K_4$-free graph with chromatic number $\chi(G) \leq 3$,
  $m$ edges, and $G \neq K_3$.  Then
  \[
    \lambda_1^2(G) + \lambda_2^2(G) \;\leq\; \frac{4m}{3}.
  \]
\end{theorem}

\begin{remark}
  Theorem~\ref{thm:3colorable} is not new: it follows immediately from
  Lin--Ning--Wu~\cite{LinNingWu2021} when $\omega(G) \leq 2$ and from
  Bollob\'as--Nikiforov~\cite{BollobasNikiforov2007} for weakly perfect
  graphs when $\omega(G) = \chi(G) = 3$.  We include it for completeness
  and as a unified summary of the $\chi \leq 3$ case.
\end{remark}

\begin{remark}
  Bollob\'as and Nikiforov~\cite{BollobasNikiforov2007} and
  Zhang (see~\cite{ElphickWocjan2014}) have proved Conjecture~\ref{conj:BN}
  for all \emph{regular} graphs.  Hence the $K_4$-free regular case is already
  resolved.  We record the precise statement for regular $K_4$-free graphs
  as a corollary of Zhang's result.
\end{remark}

\begin{corollary}[$K_4$-free regular graphs~\cite{ElphickWocjan2014}]
  \label{cor:regular}
  Let $G$ be a $d$-regular $K_4$-free graph on $n \geq 4$ vertices,
  $G \neq K_3$.  Then
  \[
    \lambda_1^2(G) + \lambda_2^2(G) \;\leq\; \frac{4m}{3}.
  \]
\end{corollary}

\begin{conjecture}[$K_4$-free graphs with large independence number]
  \label{conj:independence}
  Let $G$ be a $K_4$-free graph on $n$ vertices and $m$ edges with
  independence number $\alpha(G) \geq n/3$ and $G \neq K_3$.  Then
  \[
    \lambda_1^2(G) + \lambda_2^2(G) \;\leq\; \frac{4m}{3}.
  \]
\end{conjecture}

We present partial progress towards Conjecture~\ref{conj:independence},
identifying the exact obstruction that prevents the Hoffman-bound approach
from closing the argument.
We also formulate a spectral stability conjecture
(Conjecture~\ref{conj:stability}) for connected $K_4$-free graphs and show
that the connectivity hypothesis is necessary.

\paragraph{Proof strategy.}
Theorem~\ref{thm:multipartite} is proved in Section~\ref{sec:multipartite}
via a secular-equation analysis of the spectrum of $K_{n_1,\ldots,n_r}$;
the key step is showing that $\lambda_2(G) = 0$ whenever $n > r$.
Theorem~\ref{thm:3colorable} follows from existing results
in~\cite{BollobasNikiforov2007,LinNingWu2021} and is proved in
Section~\ref{sec:K4free}.
Corollary~\ref{cor:regular} follows from Zhang's theorem for regular graphs.
Conjecture~\ref{conj:stability} and its counterexample are presented in
Section~\ref{sec:K4free}, together with partial progress on
Conjecture~\ref{conj:independence} via the Hoffman bound~\cite{Hoffman1970}.
Section~\ref{sec:equality} characterizes the equality cases.
Section~\ref{sec:remarks} collects open problems.

%% ============================================================
\section{Preliminaries}
\label{sec:prelim}
%% ============================================================

\paragraph{Graph notation.}
All graphs are simple and undirected.
For a graph $G$ on vertex set $V(G)$ with $|V(G)| = n$ vertices and
$|E(G)| = m$ edges, write $N(v)$ for the open neighbourhood of $v$,
$d(v) = |N(v)|$ for its degree, and $\Delta(G)$, $\delta(G)$,
$\bar{d}(G) = 2m/n$ for the maximum, minimum, and average degrees.
The \emph{clique number} $\omega(G)$ is the order of the largest complete
subgraph; $\chi(G)$ is the chromatic number; $\alpha(G)$ is the
independence number.

The \emph{complete $r$-partite graph} $K_{n_1,\ldots,n_r}$ has vertex
set partitioned into $r$ independent sets (parts) of sizes
$n_1,\ldots,n_r$, with every two vertices in different parts adjacent.
Its clique number is $r$.
The \emph{Tur\'an graph} $T(n,r)$ is the unique balanced complete
$r$-partite graph on $n$ vertices, with parts of sizes
$\lfloor n/r \rfloor$ or $\lceil n/r \rceil$; it has
$e(T(n,r)) = (1-1/r)n^2/2 + O(n)$ edges.

\paragraph{Eigenvalues.}
The \emph{adjacency matrix} $A(G)$ is the symmetric $\{0,1\}$-matrix
indexed by $V(G)$ with $A_{uv} = 1$ iff $uv \in E(G)$.
Its eigenvalues are real and labelled
$\lambda_1(G) \geq \lambda_2(G) \geq \cdots \geq \lambda_n(G)$.
Since $A$ is a symmetric matrix with zero diagonal, its trace identities give
\begin{equation}
  \label{eq:trace}
  \sum_{i=1}^n \lambda_i = 0,
  \qquad
  \sum_{i=1}^n \lambda_i^2 = \tr(A^2) = 2m.
\end{equation}

The following results are used in the proofs.

\begin{theorem}[Nikiforov~\cite{Nikiforov2002}]
  \label{thm:nikiforov}
  For any graph $G$ with $m$ edges and $\omega(G) \geq 2$,
  \[
    \lambda_1(G)
    \;\leq\;
    \sqrt{2\!\left(1 - \frac{1}{\omega(G)}\right) m},
  \]
  with equality if and only if $G = T(n,\omega(G))$ for some $n$
  divisible by $\omega(G)$.
\end{theorem}

\begin{theorem}[Weyl's inequality]
  \label{thm:weyl}
  Let $A$ and $E$ be real symmetric $n \times n$ matrices.  Then
  $|\lambda_k(A + E) - \lambda_k(A)| \leq \|E\|_2$ for every $k$,
  where $\|E\|_2$ is the spectral norm of $E$.
\end{theorem}

\begin{theorem}[Hoffman bound~\cite{Hoffman1970}]
  \label{thm:hoffman}
  For any graph $G$ on $n$ vertices,
  \[
    \alpha(G)
    \;\leq\;
    \frac{-n\,\lambda_n(G)}{\lambda_1(G) - \lambda_n(G)}.
  \]
\end{theorem}

\begin{theorem}[Motzkin--Straus~\cite{MotzkinStraus1965}]
  \label{thm:MS}
  For any graph $G$ with clique number $\omega$,
  \[
    \max_{\mathbf{x} \in \Delta_n}
    \mathbf{x}^\top A(G)\, \mathbf{x}
    \;=\;
    1 - \frac{1}{\omega},
  \]
  where $\Delta_n = \{\mathbf{x} \geq 0 : \mathbf{1}^\top \mathbf{x} = 1\}$
  is the probability simplex.
  A maximizer is the uniform distribution over the vertices of any maximum
  clique.
\end{theorem}

\begin{theorem}[Interlacing theorem~\cite{HaemersInterlacing}]
  \label{thm:interlace}
  If $H$ is an induced subgraph of $G$ on $k$ vertices with eigenvalues
  $\mu_1 \geq \cdots \geq \mu_k$, then
  $\lambda_{n-k+i}(G) \leq \mu_i \leq \lambda_i(G)$ for all $i = 1, \ldots, k$.
\end{theorem}

\begin{lemma}[Equitable partition]
  \label{lem:equitable}
  Let $G$ have an equitable partition $\Pi = (V_1,\ldots,V_r)$ with quotient
  matrix $B \in \R^{r\times r}$.
  Then every eigenvalue of $B$ is an eigenvalue of $A(G)$.
  In particular $\lambda_1(G) = \lambda_1(B)$.
\end{lemma}

%% ============================================================
\section{Complete Multipartite Graphs}
\label{sec:multipartite}
%% ============================================================

We prove Theorem~\ref{thm:multipartite} by first determining the full
spectrum of $G = K_{n_1,\ldots,n_r}$.
The spectrum splits into a large zero eigenspace and $r$ further
eigenvalues determined by a secular equation.

\subsection{The spectrum of $K_{n_1,\ldots,n_r}$}

\begin{lemma}[Zero eigenspace]
  \label{lem:within}
  Let $G = K_{n_1,\ldots,n_r}$ with parts $V_1,\ldots,V_r$ of sizes
  $n_1,\ldots,n_r$.  The eigenvalue $0$ of $A(G)$ has multiplicity
  exactly $n - r$.
\end{lemma}

\begin{proof}
For each part $V_i = \{u_1,\ldots,u_{n_i}\}$ and each index
$k \in \{1,\ldots,n_i - 1\}$, define the vector $\mathbf{f}^{(i,k)}$
in $\R^n$ by
\[
  f^{(i,k)}_w
  \;=\;
  \begin{cases}
    +1 & \text{if } w = u_k, \\
    -1 & \text{if } w = u_{n_i}, \\
    \phantom{+}0 & \text{otherwise.}
  \end{cases}
\]
We claim $A(G)\mathbf{f}^{(i,k)} = \mathbf{0}$.
Take any vertex $w \in V(G)$.
If $w \in V_j$ for some $j \neq i$, then every neighbour of $w$ in
$K_{n_1,\ldots,n_r}$ lies in $V(G) \setminus V_j$, hence in particular
all neighbours of $w$ with nonzero $\mathbf{f}^{(i,k)}$-entry lie in
$V_i$.  Since $\mathbf{f}^{(i,k)}$ is supported on exactly $u_k$ and
$u_{n_i}$, both in $V_i \subseteq V(G) \setminus V_j$, we get
$(A\mathbf{f}^{(i,k)})_w = f^{(i,k)}_{u_k} + f^{(i,k)}_{u_{n_i}} = 1 + (-1) = 0$.
If $w \in V_i$, then every neighbour of $w$ lies in $V(G) \setminus V_i$,
on which $\mathbf{f}^{(i,k)}$ vanishes, so $(A\mathbf{f}^{(i,k)})_w = 0$.
Thus $\mathbf{f}^{(i,k)}$ is a $0$-eigenvector.

For fixed $i$, the vectors $\mathbf{f}^{(i,1)},\ldots,\mathbf{f}^{(i,n_i-1)}$
span the $(n_i-1)$-dimensional subspace of vectors supported on $V_i$
that sum to zero on $V_i$.
Vectors supported on different parts have disjoint supports, so the
subspaces for distinct $i$ are mutually orthogonal.
The total dimension of the zero eigenspace is at least
$\sum_{i=1}^r (n_i-1) = n - r$; the complementary $r$-dimensional
space is accounted for by the secular equation below, so the
multiplicity is exactly $n - r$.
\end{proof}

\begin{lemma}[Secular equation and the sign of $\lambda_2$]
  \label{lem:secular}
  Let $G = K_{n_1,\ldots,n_r}$ with $r \geq 2$ and parts of sizes
  $n_1 \geq \cdots \geq n_r \geq 1$.
  The eigenvalues of $A(G)$ outside the zero eigenspace are the
  real roots of
  \begin{equation}
    \label{eq:secular}
    \sum_{i=1}^r \frac{n_i}{\lambda + n_i} \;=\; 1.
  \end{equation}
  There is exactly one positive root $\alpha_1$, and all remaining
  roots are at most $-n_r < 0$.  In particular, if $n > r$ then
  $\lambda_2(G) = 0$.
\end{lemma}

\begin{proof}
\textbf{Reduction to part-constant eigenvectors.}
Any permutation of vertices within a fixed part $V_i$ is a graph
automorphism of $K_{n_1,\ldots,n_r}$, so two vertices in the same
part have identical rows in $A(G)$.
Let $\mathbf{v}$ be an eigenvector with eigenvalue $\lambda$ that is
orthogonal to the entire zero eigenspace identified in
Lemma~\ref{lem:within}.
For each $i$, the vector $\mathbf{v}$ must be constant on $V_i$:
if it were not, the projection of $\mathbf{v}$ onto the zero eigenspace
for part $V_i$ (namely the component of $\mathbf{v}$ that is supported
on $V_i$ and sums to zero there) would be nonzero, contradicting
orthogonality.
Hence we may write $v_u = c_i$ for all $u \in V_i$.

\textbf{Deriving the secular equation.}
For $u \in V_i$, the eigenvalue equation $(A\mathbf{v})_u = \lambda c_i$
reads $\sum_{j \neq i} n_j c_j = \lambda c_i$.
Setting $s := \sum_{j=1}^r n_j c_j$, this becomes $s - n_i c_i = \lambda c_i$,
so $c_i(\lambda + n_i) = s$.
For a nontrivial eigenvector at least one $c_i$ is nonzero; if $\lambda = -n_i$
for every $i$ with $c_i \neq 0$, then $s = c_i(\lambda + n_i) = 0$ for
all $i$, forcing $s = 0$ and therefore all $c_j = s/(\lambda+n_j) = 0$
for $j$ with $\lambda \neq -n_j$; the only possibility is that $\lambda = -n_i$
for all parts $i$ with $n_i$ equal to the same value, and the eigenvectors
are differences of the constant vectors on those parts --- these are already
accounted for in the zero eigenspace when all $n_i = n_j$ (since then
$\lambda + n_j = 0$, and such a difference vector has $s = 0$).
We therefore focus on $\lambda \notin \{-n_1,\ldots,-n_r\}$ and $s \neq 0$,
giving $c_i = s/(\lambda + n_i)$.
Substituting into the definition of $s$:
\[
  s
  \;=\; \sum_{i=1}^r n_i c_i
  \;=\; s \sum_{i=1}^r \frac{n_i}{\lambda + n_i},
\]
and dividing by $s \neq 0$ yields equation~\eqref{eq:secular}.

\textbf{Root analysis.}
Define $f(\lambda) = \sum_{i=1}^r n_i/(\lambda + n_i)$.
The function $f$ is continuous and strictly decreasing on each interval
between consecutive poles, since
$f'(\lambda) = -\sum_{i=1}^r n_i/(\lambda + n_i)^2 < 0$ wherever $f$ is
defined.
The poles of $f$ occur at the values $\{-n_i : 1 \leq i \leq r\}$;
let the distinct values in this set be
$-p_1 < -p_2 < \cdots < -p_s \leq -1 < 0$,
where $p_1 > p_2 > \cdots > p_s \geq 1$ are the distinct part sizes
and $s \leq r$.

\textit{One positive root.}
On $(0,+\infty)$, every term $n_i/(\lambda + n_i)$ is positive and
decreasing to $0$, so $f$ decreases strictly from $f(0) = r \geq 2 > 1$
to $\lim_{\lambda \to +\infty} f(\lambda) = 0 < 1$.
The intermediate value theorem gives a unique root $\alpha_1 \in (0,+\infty)$.

\textit{No root in $(-p_s, 0)$.}
On the interval $(-p_s, 0)$, every denominator $\lambda + n_i$ satisfies
$\lambda + n_i \geq \lambda + p_s > 0$, so $f(\lambda) \geq 0$.
Moreover, as $\lambda \to (-p_s)^+$, the term $p_s/(\lambda + p_s)$
(summed over all parts with $n_i = p_s$) diverges to $+\infty$, so
$\lim_{\lambda \to (-p_s)^+} f(\lambda) = +\infty$.
Since $f$ is strictly decreasing on $(-p_s,0)$ and $f(0) = r > 1$,
we conclude $f(\lambda) > r > 1$ for all $\lambda \in (-p_s, 0)$,
hence no root lies in this interval.

\textit{All remaining roots are at most $-p_s < 0$.}
On each inter-pole interval $(-p_{k+1},-p_k)$ for $k = 1,\ldots,s-1$:
as $\lambda \to (-p_{k+1})^+$, the terms with $n_i = p_{k+1}$ give
$f \to +\infty$, and as $\lambda \to (-p_k)^-$, the terms with
$n_i = p_k$ give $f \to -\infty$.
By the intermediate value theorem and strict monotonicity, there is
exactly one root in each such interval.
No root lies in $(-\infty, -p_1)$ because $f(\lambda) < 0 < 1$ there
(all denominators are negative and all numerators positive).
Including the poles themselves: if $\lambda = -n_j$ for some $j$ with
$\lambda + n_i \neq 0$ for all $i \neq j$, then $f(\lambda)$ is undefined,
so $\lambda$ is not a root of \eqref{eq:secular}.

\textbf{Conclusion.}
All roots of \eqref{eq:secular} outside the zero eigenspace are:
the unique positive root $\alpha_1$, and roots at most $-p_s \leq -1 < 0$.
Together with the zero eigenspace of dimension $n - r \geq 1$
(when $n > r$), the second largest eigenvalue of $A(G)$ is
$\lambda_2(G) = 0$.
\end{proof}

\begin{remark}
  \label{rem:kr}
  When $G = K_r$ (every part has size $1$), Lemma~\ref{lem:within}
  gives a zero eigenspace of dimension $n - r = 0$, and the secular
  equation $\sum_{i=1}^r 1/(\lambda+1) = 1$ has the unique positive
  solution $\lambda = r-1$ and the pole $\lambda = -1$ accounts for
  $r-1$ further eigenvalues.
  One checks $\lambda_1^2 + \lambda_2^2 = (r-1)^2 + 1 > (r-1)^2
  = 2(1-1/r)\binom{r}{2}$, which shows why $K_n$ must be excluded
  from Conjecture~\ref{conj:BN}.
\end{remark}

\subsection{Proof of Theorem~\ref{thm:multipartite}}

\begin{proof}[Proof of Theorem~\ref{thm:multipartite}]
Let $G = K_{n_1,\ldots,n_r}$ with $n \geq r+1$.
By Lemma~\ref{lem:within}, the eigenvalue $0$ has multiplicity $n-r \geq 1$.
By Lemma~\ref{lem:secular}, all eigenvalues outside the zero eigenspace
are either the unique positive value $\alpha_1$ or are strictly negative.
Ordering the full spectrum, the largest eigenvalue is $\lambda_1(G) = \alpha_1 > 0$
and the second largest is $\lambda_2(G) = 0$.

Since $\lambda_2(G) = 0$, we have $\lambda_1^2(G) + \lambda_2^2(G) = \lambda_1^2(G)$.
The clique number of $K_{n_1,\ldots,n_r}$ is $r$, so
Theorem~\ref{thm:nikiforov} gives
$\lambda_1^2(G) \leq 2(1-1/r)m$,
establishing the upper bound.

\medskip
\noindent\textbf{Equality: case $r = 2$.}
For $G = K_{n_1,n_2}$ with $r = 2$ and $n_1 + n_2 \geq 3$, we have
$\lambda_1 = \sqrt{n_1 n_2}$, $\lambda_2 = 0$, and $m = n_1 n_2$.  Hence
\[
  \lambda_1^2 + \lambda_2^2
  = n_1 n_2
  = m
  = 2\!\left(1 - \tfrac{1}{2}\right)m,
\]
so equality holds for \emph{every} complete bipartite graph $K_{n_1,n_2}$
with $n_1 + n_2 \geq 3$, regardless of whether $n_1 = n_2$.

\medskip
\noindent\textbf{Equality: case $r \geq 3$.}
Since $\lambda_2(G) = 0$, equality becomes $\lambda_1^2(G) = 2(1-1/r)m$,
which by Theorem~\ref{thm:nikiforov} holds if and only if $G = T(n,r)$,
requiring $n_1 = \cdots = n_r = n/r$.
Conversely, for the balanced partition $n_i = p$ for all $i$
(so $n = rp$ and $m = r(r-1)p^2/2$), the Perron eigenvector of the
complete $r$-partite graph assigns equal weight to each vertex, and
its Rayleigh quotient gives $\lambda_1(G) = (r-1)p$.
Therefore
\[
  \lambda_1^2(G) + \lambda_2^2(G)
  \;=\; (r-1)^2 p^2
  \;=\; \frac{2(r-1)}{r} \cdot \frac{r(r-1)p^2}{2}
  \;=\; 2\!\left(1-\frac{1}{r}\right)m,
\]
confirming equality.
For any unbalanced partition ($r \geq 3$, not all $n_i$ equal), the
positive root $\alpha_1$ of the secular equation satisfies
$\alpha_1^2 < 2(1-1/r)m$ strictly, by the equality characterization in
Theorem~\ref{thm:nikiforov}.
\end{proof}

\begin{remark}[Why $r=2$ is special]
  \label{rem:r2-equality}
  The algebraic reason every complete bipartite graph $K_{a,b}$ achieves
  equality is the identity $\lambda_1(K_{a,b}) = \sqrt{ab} = \sqrt{m}$
  combined with $\lambda_2(K_{a,b}) = 0$ (for $a + b \geq 3$).
  This forces $\lambda_1^2 + \lambda_2^2 = m = 2(1-1/2)m$ identically,
  regardless of whether $a = b$.
  For $r \geq 3$, the spectral radius $\lambda_1(K_{n_1,\ldots,n_r})$
  is strictly below $\sqrt{2(1-1/r)m}$ whenever the partition is not
  balanced, so no such universal equality holds.
\end{remark}

%% ============================================================
\section{The $K_4$-Free Case}
\label{sec:K4free}
%% ============================================================

Throughout this section $G$ is a $K_4$-free graph, so $\omega(G) \leq 3$
and the Bollob\'as--Nikiforov bound becomes $\lambda_1^2 + \lambda_2^2 \leq 4m/3$.
We write $\mathcal{T}_3$ for the family of complete tripartite graphs on $n$
vertices, and $\dedit(G, \mathcal{T}_3)$ for the minimum number of edge
insertions and deletions needed to transform $G$ into a member of $\mathcal{T}_3$.

Section~\ref{ssec:zykov} develops the Zykov symmetrization framework.
Section~\ref{ssec:3colorable} proves Theorem~\ref{thm:3colorable} and
Corollary~\ref{cor:regular}.
Section~\ref{ssec:stability} presents the open stability conjecture
(Conjecture~\ref{conj:stability}) and its counterexample.
Section~\ref{ssec:independence} presents partial progress on
Conjecture~\ref{conj:independence}.

\subsection{Zykov symmetrization}
\label{ssec:zykov}

The classical Zykov symmetrization~\cite{Zykov1949} is a graph operation
that increases (or preserves) the spectral radius while not increasing the
clique number.

\begin{definition}[Zykov operation~\cite{Zykov1949}]
  \label{def:zykov}
  Let $G$ be a graph and $u, v \in V(G)$ two non-adjacent vertices.
  The \emph{Zykov operation} $Z(G; u,v)$ replaces $u$ by a clone of $v$:
  the resulting graph $G'$ has $V(G') = V(G)$ and
  $N_{G'}(u) = N_{G'}(v) = N_G(v)$
  while all other adjacencies are unchanged.
\end{definition}

\begin{lemma}[Zykov preserves $K_4$-freeness]
  \label{lem:zykov-kfree}
  If $G$ is $K_4$-free and $u,v$ are non-adjacent, then $Z(G;u,v)$ is also
  $K_4$-free.
\end{lemma}

\begin{proof}
Any clique in $Z(G;u,v)$ containing $u$ can be transformed into a clique
in $G$ by replacing $u$ with $v$ (since $N_{Z(G;u,v)}(u) = N_G(v)$, vertex $u$
in $Z(G;u,v)$ behaves identically to $v$).
Hence $\omega(Z(G;u,v)) \leq \omega(G) \leq 3$.
\end{proof}

\begin{lemma}[Zykov increases spectral radius]
  \label{lem:zykov-lambda1}
  $\lambda_1(Z(G;u,v)) \geq \lambda_1(G)$.
\end{lemma}

\begin{proof}
Let $\mathbf{x}$ be a Perron eigenvector of $G$ with $\mathbf{x} \geq 0$.
Define $\mathbf{x}'$ by $x'_u = x'_v = \max(x_u, x_v)$ and $x'_w = x_w$
for $w \neq u, v$.
One checks that
$(\mathbf{x}')^\top A(Z(G;u,v)) \mathbf{x}' \geq \mathbf{x}^\top A(G)\mathbf{x}$,
since $u$ now has the same (or larger) neighbour-weighted sum as $v$.
By the Rayleigh quotient characterization,
$\lambda_1(Z(G;u,v)) \geq \lambda_1(G)$.
\end{proof}

\begin{remark}[Zykov and the BN ratio]
  \label{rem:zykov-ratio}
  A natural strategy would be to show that the ratio
  $(\lambda_1^2 + \lambda_2^2)/m$ is non-decreasing under each Zykov step.
  If true, since Zykov symmetrization of a $3$-colorable $K_4$-free graph
  converges to a complete tripartite graph $K_{a,b,c}$ (for which the ratio
  equals $\lambda_1^2(K_{a,b,c})/m \leq 4/3$ by
  Theorem~\ref{thm:multipartite}), one would obtain BN for the original
  graph.
  However, proving monotonicity of this ratio under a single Zykov step is
  non-trivial: both $\lambda_1^2 + \lambda_2^2$ and $m$ increase, and their
  relative rates are not easily controlled via the Rayleigh quotient alone.
  We therefore prove Theorem~\ref{thm:3colorable} via the existing
  literature in Section~\ref{ssec:3colorable} rather than via this route.
\end{remark}

\subsection{$K_4$-free graphs with $\chi(G) \leq 3$ and regular graphs}
\label{ssec:3colorable}

\begin{proof}[Proof of Theorem~\ref{thm:3colorable}]
Since $G$ is $K_4$-free, its clique number satisfies $\omega(G) \leq 3$.
We split into two cases.

\textbf{Case 1: $\omega(G) \leq 2$ (triangle-free).}
Lin, Ning, and Wu~\cite{LinNingWu2021} proved the Bollob\'as--Nikiforov
conjecture for all triangle-free graphs.  In particular,
$\lambda_1^2(G) + \lambda_2^2(G) \leq 2(1-1/2)m = m \leq 4m/3$.

\textbf{Case 2: $\omega(G) = 3$.}
From $\chi(G) \geq \omega(G) = 3$ and the hypothesis $\chi(G) \leq 3$,
we conclude $\chi(G) = \omega(G) = 3$.  A graph satisfying $\chi(G) =
\omega(G)$ is called \emph{weakly perfect}.  Bollob\'as and
Nikiforov~\cite{BollobasNikiforov2007} proved the conjecture for all
weakly perfect graphs; in particular, for $\omega = 3$:
$\lambda_1^2(G) + \lambda_2^2(G) \leq 2(1-1/3)m = 4m/3$.
\end{proof}

\begin{remark}
  The Zykov symmetrization framework of Section~\ref{ssec:zykov} provides
  an alternative proof sketch for Theorem~\ref{thm:3colorable}: iterating
  Zykov operations on a $3$-colorable $K_4$-free graph $G$ eventually
  produces a complete tripartite graph $G^*$ to which
  Theorem~\ref{thm:multipartite} applies.
  However, completing this argument requires proving that the quantity
  $(\lambda_1^2+\lambda_2^2)/m$ is non-decreasing under each Zykov step
  (see Remark~\ref{rem:zykov-ratio}); we therefore rely on the existing
  literature above.
\end{remark}

Corollary~\ref{cor:regular} is immediate from the general result of Zhang
(see~\cite{ElphickWocjan2014}), which establishes the Bollob\'as--Nikiforov
conjecture for all regular graphs regardless of clique number.

\begin{proof}[Proof of Corollary~\ref{cor:regular}]
For a $d$-regular graph $G$, $\lambda_1(G) = d$ (the all-ones vector is an
eigenvector) and $m = nd/2$.  The BN bound becomes $d^2 + \lambda_2^2 \leq
2nd/3$.
By Tur\'an's theorem, a $K_4$-free $d$-regular graph satisfies $d \leq 2n/3$,
so the right-hand side is non-negative.
The bound $d^2 + \lambda_2^2 \leq 2nd/3$ is a special case of the result
proved in~\cite{ElphickWocjan2014} for all regular graphs.
\end{proof}

\begin{remark}[Towards an independent proof for regular graphs]
  \label{rem:regular-trace}
  The following argument applies exclusively to \emph{$d$-regular} $K_4$-free
  graphs; it does not extend to irregular graphs.
  For a $d$-regular $K_4$-free graph, $K_4$-freeness implies that every
  neighbourhood $N(v)$ is triangle-free (otherwise $\{v\} \cup \triangle$
  would form a $K_4$), so Tur\'an's theorem applied to the triangle-free
  graph $G[N(v)]$ gives $e(G[N(v)]) \leq d^2/4$.
  Summing over all vertices: $3t_3(G) = \sum_v e(G[N(v)]) \leq nd^2/4$,
  so $t_3(G) \leq nd^2/12$.
  Equivalently, $\tr(A^3) = 6t_3(G) \leq nd^2/2$.
  Since $\lambda_1 = d$ for a $d$-regular graph, the trace method gives:
  \[
    \sum_{i \geq 2} \lambda_i^3 \leq \frac{nd^2}{2} - d^3 = d^2\!\left(\frac{n}{2} - d\right).
  \]
  Combined with $\sum_{i\geq 2}\lambda_i = -d$ and
  $\sum_{i\geq 2}\lambda_i^2 = nd - d^2$, a Lagrange-multiplier
  optimization of $\lambda_2$ subject to these three moment constraints yields
  an upper bound on $\lambda_2$.  Closing this bound to $\sqrt{d(2n/3-d)}$
  via the trace method alone appears to require the fourth moment
  $\tr(A^4)$, which in turn demands a sharp estimate on $\sum_{uv\in E}|N(u)\cap
  N(v)|^2$ for $K_4$-free graphs---a quantity tightly related to the number of
  $4$-cycles.  We leave the completion of this direct argument as an open
  problem.
\end{remark}

\subsection{Spectral stability for dense $K_4$-free graphs (open)}
\label{ssec:stability}

A natural approach to the full conjecture for dense $K_4$-free graphs is
via a stability argument: if $\lambda_1^2 + \lambda_2^2$ is close to the
maximum $4m/3$, the graph should be close to the balanced tripartite graph
$T(n,3)$, and one could hope to bootstrap this structural proximity into the
exact bound.
However, as the counterexample in Remark~\ref{rem:stability-gap} shows,
such an argument cannot succeed without a connectivity hypothesis,
and we are not aware of a reference that supplies a spectral stability lemma
in the required generality.
Note also that a Weyl-inequality perturbation argument applied to a nearby
tripartite graph yields only the asymptotic bound
$\lambda_1^2 + \lambda_2^2 \leq \tfrac{4m}{3} + O(\varepsilon m)$,
not the exact inequality.
We therefore state the following as an open problem.

\begin{conjecture}[Spectral stability for $K_4$-free graphs]
  \label{conj:stability}
  For every $\varepsilon > 0$ there exists $\delta > 0$ such that
  the following holds.
  Let $G$ be a $K_4$-free \emph{connected} graph on $n$ vertices and
  $m$ edges with
  $\lambda_1^2(G) + \lambda_2^2(G) > (4/3 - \delta)m$.
  Then $\dedit(G, \mathcal{T}_3) \leq \varepsilon n^2$.
\end{conjecture}

\begin{remark}[Why the connectivity hypothesis is necessary]
  \label{rem:stability-gap}
  Without the connectivity assumption, Conjecture~\ref{conj:stability}
  is false.
  As a counterexample, let $G = K_{k,k,k} \cup kK_1$ (the balanced
  complete tripartite graph on $3k$ vertices together with $k$ isolated
  vertices).
  Then $n = 4k$, $m = 3k^2$, $\omega(G) = 3$, $\lambda_1 = 2k$,
  $\lambda_2 = 0$, and
  \[
    \lambda_1^2 + \lambda_2^2 = 4k^2 = \tfrac{4}{3}m,
  \]
  so the spectral hypothesis holds for every $\delta > 0$.
  However, any complete tripartite graph on $4k$ vertices has approximately
  $\tfrac{2}{3}(4k)^2/2 = \tfrac{16k^2}{3} \approx 5.3k^2$ edges, whereas
  $G$ has only $3k^2$, so $\dedit(G, \mathcal{T}_3) = \Omega(k^2) = \Omega(n^2)$.

  Moreover, applying Weyl's inequality in the stability argument yields
  only $\lambda_1^2 + \lambda_2^2 \leq \tfrac{4m}{3} + O(\varepsilon m)$,
  not the exact bound $\leq \tfrac{4m}{3}$.
  Proving the sharp form Conjecture~\ref{conj:stability} for dense
  connected $K_4$-free graphs therefore requires additional structural
  information beyond what spectral stability alone provides.
\end{remark}

\subsection{Partial results for $K_4$-free graphs with $\alpha(G) \geq n/3$}
\label{ssec:independence}

We identify a natural approach to Conjecture~\ref{conj:independence}
and determine precisely why it falls short.

\begin{proposition}
  \label{prop:hoffman}
  Let $G$ be a $K_4$-free graph on $n$ vertices and $m$ edges with
  $\alpha(G) \geq n/3$.  Then $|\lambda_n(G)| \geq \lambda_1(G)/2$.
\end{proposition}

\begin{proof}
The Hoffman bound (Theorem~\ref{thm:hoffman}) gives
$\alpha(G) \leq -n\lambda_n/(\lambda_1 - \lambda_n)$.
Write $\mu = |\lambda_n(G)| = -\lambda_n(G) \geq 0$.
Substituting $\alpha(G) \geq n/3$:
\[
  \frac{n}{3}
  \;\leq\; \frac{n\mu}{\lambda_1 + \mu}.
\]
Cross-multiplying by $3(\lambda_1 + \mu) > 0$ gives
$\lambda_1 + \mu \leq 3\mu$, i.e., $\lambda_1 \leq 2\mu$.
\end{proof}

\begin{proposition}[Obstruction to closing the argument]
  \label{prop:obstruction}
  Let $G$ be a $K_4$-free graph on $n$ vertices and $m$ edges with
  $\alpha(G) \geq n/3$ and $G \neq K_3$.
  The bound $|\lambda_n(G)| \geq \lambda_1(G)/2$ from
  Proposition~\ref{prop:hoffman} and the trace identity together give
  \[
    \lambda_1^2(G) + \lambda_2^2(G)
    \;\leq\; 2m - \frac{\lambda_1^2(G)}{4}.
  \]
  This bound is at most $4m/3$ if and only if $\lambda_1^2(G) \geq 8m/3$,
  a condition that is never satisfied for $K_4$-free graphs.
\end{proposition}

\begin{proof}
From the trace identity $\sum_{i=1}^n \lambda_i^2 = 2m$ and
$\lambda_n^2 \geq \lambda_1^2/4$ (Proposition~\ref{prop:hoffman}):
\[
  \lambda_1^2 + \lambda_2^2
  \;\leq\; 2m - \lambda_n^2
  \;\leq\; 2m - \frac{\lambda_1^2}{4}.
\]
For this upper bound to imply $\lambda_1^2 + \lambda_2^2 \leq 4m/3$, we would
need $2m - \lambda_1^2/4 \leq 4m/3$, i.e., $\lambda_1^2 \geq 8m/3$.
However, Theorem~\ref{thm:nikiforov} with $\omega(G) \leq 3$ gives
$\lambda_1^2 \leq 4m/3 < 8m/3$.
Hence the Hoffman-energy approach is provably insufficient, and the
bound it produces, namely $2m - \lambda_1^2/4 \geq 2m - m/3 = 5m/3 > 4m/3$,
is too weak by a factor of $5/4$.
\end{proof}

The obstruction in Proposition~\ref{prop:obstruction} is not merely a
deficiency of the method: it pinpoints exactly what additional eigenvalue
information would close the argument.
Any proof of Conjecture~\ref{conj:independence} must use structural
properties of $G$ beyond the Hoffman bound and the trace identity.
For the case $\omega(G) \leq 2$ (triangle-free), Theorem~\ref{thm:3colorable}
already settles the conjecture.
For $\omega(G) = 3$ with $\alpha(G) \geq n/3$, the fractional chromatic
number satisfies $\chi_f(G) \leq n/\alpha(G) \leq 3$, placing this case
at an intermediate level of difficulty between the weakly-perfect case
(proved) and the full $K_4$-free conjecture (open).

%% ============================================================
\section{Equality in the Complete Multipartite Case}
\label{sec:equality}
%% ============================================================

\begin{theorem}[Equality characterization for proved cases]
  \label{thm:equality}
  Let $G \neq K_n$ with $m$ edges and $\omega = \omega(G)$.
  Suppose $\lambda_1^2(G) + \lambda_2^2(G) = 2(1 - 1/\omega)m$.
  \begin{enumerate}[label=(\alph*)]
    \item If $\omega = 2$ and $G$ is triangle-free, then $G$ is a complete
          bipartite graph $K_{a,b}$ with $a + b \geq 3$.
    \item If $G = K_{n_1,\ldots,n_r}$ is a complete $r$-partite graph
          with $r = \omega \geq 2$ and $n \geq r+1$, then:
          either $r = 2$ (and equality holds for every $K_{a,b}$ with
          $a+b\geq 3$), or $r \geq 3$ with $n_1 = \cdots = n_r$ (i.e.,
          $G = T(n,r)$).
  \end{enumerate}
\end{theorem}

\begin{proof}
\textbf{Part (a).}
For triangle-free $G$ with $\omega = 2$, Lin--Ning--Wu~\cite{LinNingWu2021}
proved Conjecture~\ref{conj:BN} with equality characterization:
equality holds iff $G$ is a complete bipartite graph.
(Directly: $K_{a,b}$ has $\lambda_1 = \sqrt{ab}$, $\lambda_2 = 0$,
$m = ab$, so $\lambda_1^2 + \lambda_2^2 = ab = m = 2(1-1/2)m$.)

\textbf{Part (b).}
For $G = K_{n_1,\ldots,n_r}$ with $n \geq r+1$, Lemma~\ref{lem:secular}
gives $\lambda_2(G) = 0$.  Hence $\lambda_1^2 + \lambda_2^2 = \lambda_1^2$.
Equality $\lambda_1^2 = 2(1-1/r)m$ is then exactly the equality case of
Theorem~\ref{thm:nikiforov}, and is characterized by
Theorem~\ref{thm:multipartite}: all parts equal (for $r \geq 3$) or any
complete bipartite graph (for $r = 2$).
\end{proof}

\begin{remark}
  For general graphs $G$ (not necessarily complete multipartite),
  characterizing all equality cases of Conjecture~\ref{conj:BN} requires
  a complete proof of the conjecture, which remains open.
  We therefore restrict Theorem~\ref{thm:equality} to the classes for which
  the conjecture has been proved.
  The general equality condition stated in Conjecture~\ref{conj:BN} is
  itself conjectured.
\end{remark}

%% ============================================================
\section{Remarks and Open Problems}
\label{sec:remarks}
%% ============================================================

Theorem~\ref{thm:multipartite} resolves the Bollob\'as--Nikiforov
conjecture completely for complete multipartite graphs, a family that
includes the equality cases and the Tur\'an graphs.
The case that remains open is $K_4$-free graphs with $\chi(G) \geq 4$
and $m = o(n^2)$ or small $n$; this includes Mycielski-type constructions,
which achieve high chromatic number while being $K_4$-free.

\paragraph{Summary of results.}
\begin{enumerate}
  \item \textbf{New result:} The Bollob\'as--Nikiforov conjecture holds for
        every complete multipartite graph $K_{n_1,\ldots,n_r}$ with $n > r$
        (Theorem~\ref{thm:multipartite}).
        The proof uses a secular-equation characterization of the spectrum,
        establishing $\lambda_2 = 0$ for such graphs, and then applies
        Nikiforov's spectral Tur\'an theorem.
        The equality cases are: every complete bipartite graph $K_{a,b}$
        ($r=2$, $a+b \geq 3$), and the balanced complete $r$-partite graph
        $T(n,r)$ ($r \geq 3$, equal parts).

  \item \textbf{Known results unified:} The conjecture holds for $K_4$-free
        graphs with $\chi(G) \leq 3$ (Theorem~\ref{thm:3colorable}),
        following from~\cite{LinNingWu2021} and~\cite{BollobasNikiforov2007};
        and for all regular graphs (Corollary~\ref{cor:regular}) by
        Zhang's theorem~\cite{ElphickWocjan2014}.

  \item \textbf{Stability (open):} Conjecture~\ref{conj:stability}
        proposes a spectral stability statement for connected $K_4$-free
        graphs.
        The counterexample $G = K_{k,k,k} \cup kK_1$
        (Remark~\ref{rem:stability-gap}) shows connectivity is necessary:
        $G$ achieves $\lambda_1^2 + \lambda_2^2 = 4m/3$ yet requires
        $\Omega(n^2)$ edge edits to become tripartite.
        We are not aware of a reference that supplies a spectral stability
        theorem in the generality required to prove the conjecture.

  \item \textbf{Partial progress:} For $K_4$-free graphs with
        $\alpha(G) \geq n/3$, the Hoffman bound gives
        $|\lambda_n| \geq \lambda_1/2$, which implies $\lambda_1^2 \leq 8m/5$.
        The obstruction (Proposition~\ref{prop:obstruction}) shows that
        any proof of Conjecture~\ref{conj:independence} must go beyond
        the Hoffman bound and the trace identity.
\end{enumerate}

\paragraph{Open problems.}
\begin{enumerate}
  \item Prove Conjecture~\ref{conj:independence}: BN for $K_4$-free graphs
        with $\alpha(G) \geq n/3$.
        The spectral obstruction identified in
        Section~\ref{ssec:independence} suggests that a new eigenvalue
        lower bound---beyond what Hoffman and Nikiforov provide---is needed.

  \item Prove Conjecture~\ref{conj:stability}: sharp BN stability for
        dense connected $K_4$-free graphs.
        The counterexample $K_{k,k,k} \cup kK_1$ shows connectivity is
        necessary.

  \item \textbf{(Main open problem)} Prove the Bollob\'as--Nikiforov
        conjecture for all $K_4$-free graphs with $\chi(G) \geq 4$.

  \item Prove the conjecture for $K_5$-free graphs ($\omega = 4$).
        Theorem~\ref{thm:multipartite} covers complete $4$-partite graphs;
        the general $K_5$-free case requires new ideas.

  \item Determine, for each $\omega \geq 3$, the best constant $C_\omega$
        such that $\lambda_1^2 + \lambda_2^2 \leq C_\omega \cdot m$ for all
        $K_{\omega+1}$-free graphs with $\chi \geq \omega+1$.

  \item Complete the direct trace-method proof of BN for regular
        $K_4$-free graphs (Remark~\ref{rem:regular-trace}).
\end{enumerate}

%% ============================================================
%%  Acknowledgements
%% ============================================================
\section*{Acknowledgements}
The author thanks Dr. Shengtong Zhang for helpful discussions.

%% ============================================================
%%  References
%% ============================================================
\bibliographystyle{elsarticle-num}
\begin{thebibliography}{99}

\bibitem{BollobasNikiforov2007}
B.~Bollob\'as, V.~Nikiforov,
Cliques and the spectral radius,
\textit{J. Combin. Theory Ser.~B} \textbf{97} (2007) 859--865.

\bibitem{ElphickWocjan2014}
C.~Elphick, P.~Wocjan,
An inertial lower bound for the chromatic number of a graph,
\textit{Electron. J. Combin.} \textbf{24} (2017) \#P1.56.

\bibitem{HaemersInterlacing}
W.~Haemers,
Interlacing eigenvalues and graphs,
\textit{Linear Algebra Appl.} \textbf{226--228} (1995) 593--616.

\bibitem{Hoffman1970}
A.J.~Hoffman,
On eigenvalues and colorings of graphs,
in: \textit{Graph Theory and its Applications} (B.~Harris, Ed.),
Academic Press, 1970, pp.~79--91.

\bibitem{KumarPragada2024}
H.~Kumar, S.~Pragada,
Bollob\'as--Nikiforov conjecture for graphs with not so many triangles,
\textit{arXiv:2407.19341}, 2024.

\bibitem{LinNingWu2021}
H.~Lin, B.~Ning, B.~Wu,
Eigenvalues and triangles in graphs,
\textit{Combin. Probab. Comput.} \textbf{30} (2021) 258--270.

\bibitem{LiuBu2025}
Y.~Liu, Z.~Bu,
Bollob\'as--Nikiforov conjecture holds asymptotically almost surely,
\textit{arXiv:2501.07137}, 2025.

\bibitem{MotzkinStraus1965}
T.S.~Motzkin, E.G.~Straus,
Maxima for graphs and a new proof of a theorem of Tur\'an,
\textit{Canad. J. Math.} \textbf{17} (1965) 533--540.

\bibitem{Nikiforov2002}
V.~Nikiforov,
Some inequalities for the largest eigenvalue of a graph,
\textit{Combin. Probab. Comput.} \textbf{11} (2002) 179--189.

\bibitem{Nosal1970}
E.~Nosal,
\textit{Eigenvalues of Graphs},
Master's Thesis, University of Calgary, 1970.

\bibitem{Zykov1949}
A.A.~Zykov,
On some properties of linear complexes,
\textit{Mat. Sb.} \textbf{24} (1949) 163--188.

\end{thebibliography}

\end{document}
